% Table tab:literature_comparison_resultssec — comparison of mutation prediction
% methods in NSCLC. Synced with thesis Table 2.5 (ch02_background.tex, lines
% 1299-1331): seven prior studies + this work's two MIL conditions.
% AUROC values are means or ranges across the cited paper's CV protocol;
% daggers (^\dagger) indicate values approximated from published figures.
% Wide table: wrapped in \resizebox so it scales to the printed column width.
\resizebox{\linewidth}{!}{%
\setlength{\tabcolsep}{4pt}
\begin{tabular}{l c l r l c c c c c c}
\toprule
\textbf{Study} & \textbf{Year} & \textbf{Journal} & $\boldsymbol{N}$ &
  \textbf{Method} &
  \textit{TP53} & \textit{EGFR} & \textit{KRAS} &
  \textit{STK11} & \textit{KEAP1} & \textit{RBM10} \\
\midrule
Coudray~et~al.\cite{coudray2018classification} &
  2018 & Nat.\ Med.\ & 382 & Inception-v3 e2e\textsuperscript{*} &
  0.760 & 0.826 & 0.733 & 0.856 & 0.675 & --- \\
Kather~et~al.\cite{kather2020pancancer} &
  2020 & Nat.\ Cancer & 17{,}355 & ShuffleNet+att.\ &
  0.65--0.72$^{\dagger}$ & 0.63--0.68$^{\dagger}$ & 0.60--0.65$^{\dagger}$ & --- & --- & --- \\
Saldanha~et~al.\cite{saldanha2023npj} &
  2023 & npj Prec.\ Onc.\ & 676 & RetCCL+attMIL &
  --- & 0.76\textsuperscript{\P} & --- &
  --- & --- & --- \\
Zhao~et~al.\cite{zhao2025deepgem} (DeepGEM) &
  2025 & Lancet Onc.\ & 3{,}637 & co-supervised MIL &
  0.82--0.96\textsuperscript{$\ddagger$} & 0.945\textsuperscript{$\ddagger$} & 0.82--0.96\textsuperscript{$\ddagger$} & --- & --- & --- \\
Campanella~et~al.\cite{campanella2025eagle} (EAGLE) &
  2025 & Nat.\ Med.\ & 8{,}461 & FM fine-tuning &
  --- & 0.890 & --- & --- & --- & --- \\
Wang~et~al.\cite{wang2024chief} (CHIEF) &
  2024 & Nature & 60{,}530 & FM, pan-cancer &
  ---\textsuperscript{\S} & $>$\,0.70 & --- & --- & --- & --- \\
\midrule
\multirow{3}{*}{\textbf{This work}} & \multirow{3}{*}{2026} & \multirow{3}{*}{---} & \multirow{3}{*}{\textbf{458}}
  & B2 (encoder only)\textsuperscript{\textbardbl} & \textbf{0.661} & \textbf{0.582} & \textbf{0.517} & 0.591 & \textbf{0.488} & 0.536 \\
  & & & & \piabmil{} & 0.658 & 0.580 & 0.516 & 0.578 & 0.446 & \textbf{0.552} \\
  & & & & \fcmil{} & 0.646 & 0.567 & 0.506 & \textbf{0.608} & 0.460 & 0.528 \\
\bottomrule
\end{tabular}%
}
\begin{flushleft}\footnotesize
\textsuperscript{*}The whole Coudray row is the per-slide AUC obtained by
averaging the predicted tile probabilities (their Fig.~3 table), measured on a
$62$-slide held-out test set. The widely quoted $0.733$--$0.856$ is the range
over the six genes they call predictable, and its endpoints are \textit{KRAS}
and \textit{STK11}, not a per-gene value for either \textit{TP53} or any other
gene. Their intervals are correspondingly wide, e.g.\ \textit{STK11}
$[0.709,0.964]$ and \textit{KEAP1} $[0.466,0.865]$.
$^{\dagger}$Kather~et~al.\ report LUAD per-gene AUROCs only in a bar chart
(their Extended Data Fig.~2); the values shown are read from that figure and lie
within the paper's stated LUAD range of $0.60$--$0.78$ (3-fold CV on TCGA).
\textsuperscript{\P}Saldanha~et~al.\ report only \textit{EGFR} for LUAD
(AUROC $0.76\!\pm\!0.03$, external validation on CPTAC); the other LUAD
target genes are not individually reported in the source.
\textsuperscript{$\ddagger$}DeepGEM reports an overall cross-gene AUC range
on the public TCGA set (0.82--0.96) rather than per-gene values; EGFR is
the strongest (${\approx}0.945$) and STK11 is not among its target genes.
\textsuperscript{\S}CHIEF's validated LUAD-specific result is EGFR
($>0.70$ in held-out and independent sets); its strong TP53 signal is in
other tumour types (low-grade glioma, adrenal, endometrial), not LUAD.
EAGLE's 0.890 is the prospective silent-trial AUC (retrospective means:
internal 0.847, external 0.870).
Cohort sizes ($N$) are slide counts.
``---'' indicates the gene was not reported in the corresponding study;
Coudray is the only prior study reporting \textit{STK11} and \textit{KEAP1}
individually, and \textit{RBM10} is reported by none of them.
\textsuperscript{\textbardbl}Our rows are pooled out-of-fold \auroc{},
patient-preserved folds, averaged over five seeds
(\cref{tab:mutation_luad_resultssec}); B2 is the encoder-only ceiling
baseline and is included because it, not either fusion condition, is the
best this cohort yields on four of six genes.
\par\smallskip
\textbf{This row is not commensurable with the rows above it, and the gap
should not be read as a method comparison.} Ours are scored under nested
cross-validation, with early stopping on an inner split and the outer fold
scored exactly once, on a cohort restricted to adenocarcinoma cases with a
mutation call and cleared of solid-tissue-normal sections. The prior line
predominantly reports values selected on the same fold they are scored on,
and on inventories that retain squamous and normal sections. Removing those
two sources of optimism is what separates the two blocks; a like-for-like
comparison would require rerunning each prior method under this protocol,
which we do not do. The audit of what these cohorts can resolve at all is
\cref{tab:literature_power_audit}.
\end{flushleft}
